%%%%%%%%%%%%%%%%%%%%%%%%%%%%%%%%%%%%%%%%%%%%%%%%%%%%%%%%%%%%%%%%%%%
%                                                                 %
%                            CHAPTER ONE                          %
%                                                                 %
%%%%%%%%%%%%%%%%%%%%%%%%%%%%%%%%%%%%%%%%%%%%%%%%%%%%%%%%%%%%%%%%%%%


%\renewcommand{\simplechapterdelim}{.}
%\simplechapter

\chapter{Εισαγωγή}

Οι χρηματοοικονομικές αγορές αποτελούν ένα από τα πιο απαιτητικά πεδία εφαρμογής της τεχνητής νοημοσύνης, καθώς χαρακτηρίζονται από έντονη αβεβαιότητα, μη-στασιμότητα, μεταβολές καθεστώτων (regime shifts), θόρυβο στις παρατηρήσεις και ισχυρούς πρακτικούς περιορισμούς (κόστη συναλλαγών, ρευστότητα, περιορισμοί χαρτοφυλακίου). Σε αυτό το περιβάλλον, οι παραδοσιακές προσεγγίσεις πρόβλεψης τιμών ή σταθερών κανόνων συναλλαγών συχνά δυσκολεύονται να γενικεύσουν εκτός δείγματος, ιδίως όταν το υποκείμενο δυναμικό σύστημα αλλάζει. Παράλληλα, η εκτεταμένη παραγωγή ειδησεογραφικού περιεχομένου και σχολιασμού στα κοινωνικά δίκτυα δημιουργεί μια πρόσθετη πηγή πληροφορίας, η οποία επηρεάζει τις προσδοκίες και κατ’ επέκταση την αγορά, αλλά είναι δύσκολο να ενσωματωθεί με συνεπή και αυτοματοποιημένο τρόπο σε έναν πράκτορα συναλλαγών.

Η παρούσα εργασία, με τίτλο ''Τεχνικές βαθιάς μάθησης για υλοποίηση προσαρμοστικών πρακτόρων με εφαρμογή στις χρηματοοικονομικές συναλλαγές'', διερευνά μια ενοποιημένη προσέγγιση όπου ο πράκτορας ενισχύεται με \emph{sentiment} που παράγεται από ένα Μεγάλο Γλωσσικό Μοντέλο (LLM). Η βασική υπόθεση είναι ότι το κειμενικό σήμα (ειδήσεις, άρθρα, αναρτήσεις) μπορεί να προσφέρει συμπληρωματική πληροφορία σε σχέση με τα κλασικά αριθμητικά χαρακτηριστικά αγοράς, ιδιαίτερα σε περιόδους όπου η μεταβλητότητα, η αβεβαιότητα ή η κατεύθυνση της τάσης καθορίζονται σε σημαντικό βαθμό από την πληροφόρηση και τη συλλογική ψυχολογία. Ωστόσο, η αξιοποίηση του sentiment δεν είναι προφανής: απαιτεί αξιόπιστη εξαγωγή από κείμενο, κατάλληλη αναπαράσταση στο χώρο χαρακτηριστικών και, κυρίως, τρόπο ενσωμάτωσης στο μηχανισμό λήψης αποφάσεων ώστε να ενισχύει την προσαρμοστικότητα χωρίς να οδηγεί σε ευθραυστότητα ή υπερπροσαρμογή.

Για την επίτευξη αυτού του στόχου, ο προτεινόμενος πράκτορας οργανώνεται σε δύο βασικούς κορμούς. Ο πρώτος κορμός αφορά την παραγωγή sentiment μέσω LLM, το οποίο έχει προσαρμοστεί σε σχετικό domain με χρήση τεχνικών αποδοτικής εκπαίδευσης παραμέτρων, συγκεκριμένα LoRA, αξιοποιώντας δεδομένα από άρθρα και αναρτήσεις. Ο δεύτερος κορμός αφορά έναν πράκτορα ενισχυτικής μάθησης με βαθύ αναδρομικό νευρωνικό δίκτυο (LSTM), ο οποίος εκπαιδεύεται με Proximal Policy Optimization (PPO). Το sentiment αξιοποιείται διττά: (α) ως πρόσθετο χαρακτηριστικό εισόδου του μοντέλου πολιτικής/κριτικού, ώστε να εμπλουτίζει την κατάσταση του περιβάλλοντος με κειμενικά-σημασιολογική πληροφορία, και (β) ως μηχανισμός ''amplify/reduce'' στις ενέργειες του LSTM, δηλαδή ως σήμα που ρυθμίζει την ένταση/κατεύθυνση της απόφασης συναλλαγής υπό συγκεκριμένες συνθήκες. Με αυτόν τον τρόπο, το sentiment δεν αντιμετωπίζεται απλώς ως παθητικό feature, αλλά ως παράγοντας που μπορεί να επηρεάζει λειτουργικά την εκτελεστική συμπεριφορά του πράκτορα.

Ένα κρίσιμο ζήτημα σε τέτοιες προσεγγίσεις είναι η αξιόπιστη εκτίμηση της γενίκευσης, δεδομένου ότι οι χρηματοοικονομικές χρονοσειρές παρουσιάζουν ισχυρή χρονική εξάρτηση και οι απλές τυχαίες διασπάσεις train/test είναι συχνά παραπλανητικές. Για τον λόγο αυτό, η εργασία ενσωματώνει διαδικασία επιλογής υπερπαραμέτρων με Combinatorial Purged Cross-Validation (CPCV), η οποία είναι καταλληλότερη για χρονικά εξαρτώμενα δεδομένα, καθώς μειώνει τη διαρροή πληροφορίας (look-ahead bias) και επιτρέπει πιο συντηρητική και ρεαλιστική αξιολόγηση. Η CPCV αξιοποιείται για την αναζήτηση βέλτιστων υπερπαραμέτρων του PPO--LSTM (και των σχετικών ρυθμίσεων ενσωμάτωσης sentiment), ώστε η τελική πολιτική να προκύπτει από διαδικασία που στοχεύει στη σταθερότητα και την ανθεκτικότητα εκτός δείγματος.

Συνοψίζοντας, η συμβολή της εργασίας έγκειται στη μελέτη ενός προσαρμοστικού πράκτορα συναλλαγών που γεφυρώνει (i) μοντελοποίηση κειμενικής πληροφορίας με LLM προσαρμοσμένο μέσω LoRA και (ii) λήψη αποφάσεων σε ακολουθιακό περιβάλλον με PPO πάνω σε LSTM, με το sentiment να λειτουργεί τόσο ως πληροφοριακό χαρακτηριστικό όσο και ως σήμα ρύθμισης των ενεργειών. Η μεθοδολογία πλαισιώνεται από CPCV για την επιλογή υπερπαραμέτρων, με στόχο μια πιο αυστηρή προσέγγιση αξιολόγησης σε μη-στασιμό χρόνο.

Το υπόλοιπο της εργασίας οργανώνεται ως εξής: στο επόμενο κεφάλαιο παρουσιάζεται το θεωρητικό υπόβαθρο (ενισχυτική μάθηση, PPO, LSTM, LLM και LoRA) και οι σχετικές έννοιες εξαγωγής/χρήσης sentiment. Στη συνέχεια περιγράφεται η προτεινόμενη αρχιτεκτονική και ο τρόπος ενσωμάτωσης του sentiment ως feature και ως μηχανισμού διαμόρφωσης ενεργειών. Ακολουθεί η πειραματική μεθοδολογία, συμπεριλαμβανομένης της CPCV διαδικασίας για υπερπαραμέτρους, και τέλος παρουσιάζονται τα αποτελέσματα, η συζήτηση και τα συμπεράσματα με προτάσεις για μελλοντική εργασία.
